\begin{table}[htbp]
    \centering
    \caption{Summary of Change in Opioid Prescribing Rate by Authority and Provider Type}
    \label{tab:d_partd_opioid_rate_provtype}
    \vspace{0.2cm}
    \begin{tabular}{lccc}
        \toprule
        & \multicolumn{3}{c}{\textbf{Provider Type}} \\
        \cmidrule(lr){2-4}
        \textbf{State NP Authority} & \textbf{MD/DO} & \textbf{Nurse Practitioner} & \textbf{Physician Assistant} \\
        \midrule
        Restricted & -0.007 & -0.008 & -0.013 \\
        Reduced & -0.006 & -0.009 & -0.014 \\
        Full Practice & -0.006 & -0.010 & -0.015 \\
        \midrule
        \textbf{National Average} & \multicolumn{3}{c}{-0.007} \\
        \bottomrule
    \end{tabular}
    \vspace{0.1cm}
    \begin{minipage}{0.85\textwidth}
        \footnotesize \textit{Notes:} Outcome represents year-over-year absolute change. Baseline Metric: Rate: Proportion of total Part D claims that are Schedule II/III opioids. NP Law categories represent Restricted (supervision), Reduced (collaborative), and Full (independent practice). Categorization reflects the regulatory environment as of 2023.
    \end{minipage}
\end{table}